\documentclass[11pt,a4paper]{article}
\usepackage{fontspec,amsmath,amssymb,amsthm,microtype}
\setmainfont{lmroman10-regular.otf}[BoldFont=lmroman10-bold.otf,ItalicFont=lmroman10-italic.otf,BoldItalicFont=lmroman10-bolditalic.otf]
\usepackage{xcolor}
\usepackage[most]{tcolorbox}
\usepackage{mathrsfs}
\providecommand{\lt}{<}
\providecommand{\gt}{>}
\setmonofont{DejaVuSansMono.ttf}[BoldFont=DejaVuSansMono-Bold.ttf,ItalicFont=DejaVuSansMono-Oblique.ttf,BoldItalicFont=DejaVuSansMono-BoldOblique.ttf,Scale=0.86]
\newfontfamily\leanfallback{STIXTwoMath-Regular.otf}[Scale=0.86]
\definecolor{leanrule}{gray}{0.62}
\definecolor{leanback}{gray}{0.962}
\definecolor{leantitle}{gray}{0.22}
\definecolor{leanmuted}{gray}{0.38}
\newtcolorbox{leanbox}[1]{enhanced,breakable,colback=leanback,colframe=leanrule,
  boxrule=0.4pt,arc=0.6mm,left=2.4mm,right=2.4mm,top=1.8mm,bottom=1.8mm,
  before skip=0.8em,after skip=1em,fonttitle=\small\bfseries,coltitle=leantitle,
  colbacktitle=leanback,titlerule=0.3pt,title={#1}}
\newlength{\leancharwidth}
\newenvironment{leancode}{\par\smallskip\begingroup\ttfamily\footnotesize\raggedright
  \settowidth{\leancharwidth}{M}\setlength{\parskip}{0pt}\setlength{\parindent}{0pt}
  \frenchspacing}{\par\endgroup\smallskip}
\newcommand{\leanline}[2]{\hangindent=\dimexpr#1\leancharwidth+4\leancharwidth\relax
  \hangafter=1\noindent\hspace*{#1\leancharwidth}#2\par}
\newcommand{\leansource}[1]{\par{\footnotesize\color{leanmuted}#1}\par}
\newcommand{\leanlabel}[1]{\par\smallskip\noindent{\small\bfseries #1}\par}
\usepackage[margin=25mm]{geometry}
\usepackage[unicode,hidelinks]{hyperref}
\hypersetup{pdfauthor={LangevinLean Project}}
\setlength{\emergencystretch}{3em}
\setlength{\parskip}{0.4em}
\setlength{\parindent}{0pt}
\allowdisplaybreaks[2]
\raggedbottom
\title{\texorpdfstring{Scalar convolution and primitive identities}{Scalar convolution and primitive identities}}
\author{LangevinLean Project}
\date{}
\begin{document}\maketitle
Draft mathematical exposition; these notes may contain mistakes.


The two general coefficient identities are proved for every pair of natural indices. The primitive identity starts at \(\varphi _{1}\), as required by the definition \(\varphi _{0}=0\); the convolution identity includes both zero indices. The proofs hold for every real friction and time.



Integration by parts transfers primitive indices between convolution factors. The resulting friction and derivative recurrences give the exact correction relations, with the base-index exception stated explicitly.


\section{\texorpdfstring{Move a primitive between the two convolution factors}{Move a primitive between the two convolution factors}}


Differentiate \(\psi_{k+1}(s)\psi_{j+1}(t-s)\). Its derivative is \(\psi_k(s)\psi_{j+1}(t-s)-\psi_{k+1}(s)\psi_j(t-s)\). Both boundary products vanish because every positive-index response vanishes at zero. One proves the derivative by the product and chain rules, proves both continuous terms integrable, and applies the fundamental theorem. Therefore \[\int_0^t\psi_k(s)\psi_{j+1}(t-s)\,ds=\int_0^t\psi_{k+1}(s)\psi_j(t-s)\,ds.\] Induction moves the entire second index to the first. At second index zero the defining exponential convolution gives exactly \(\varphi_{k+1}\), proving \(\varphi_{k+j+1}=\psi_k*\psi_j\) for all \(k,j\ge0\).



The boundary terms vanish because positive-index response primitives are zero at the origin. Integration by parts can therefore transfer one primitive from one convolution factor to the other without a remainder. Repeating this operation yields the full index-addition convolution identity.

\begin{leanbox}{Lean: \texttt{psi\_\allowbreak{}convolution\_\allowbreak{}shift}, \texttt{phi\_\allowbreak{}eq\_\allowbreak{}psi\_\allowbreak{}convolution}}
\leansource{\texttt{psi\_\allowbreak{}convolution\_\allowbreak{}shift} \quad KineticLangevin/\allowbreak{}ScalarConvolution.lean, lines 8–23}
\begin{leancode}
\leanline{0}{theorem\ psi\_convolution\_shift\ (γ\ t\ :\ ℝ)\ (k\ j\ :\ ℕ)\ :}
\leanline{4}{(∫\ s\ in\ 0..t,psi\ γ\ k\ s*psi\ γ\ (j+1)\ (t-s))=}
\leanline{6}{∫\ s\ in\ 0..t,psi\ γ\ (k+1)\ s*psi\ γ\ j\ (t-s)}
\end{leancode}
\leansource{\texttt{phi\_\allowbreak{}eq\_\allowbreak{}psi\_\allowbreak{}convolution} \quad KineticLangevin/\allowbreak{}ScalarConvolution.lean, lines 26–37}
\begin{leancode}
\leanline{0}{theorem\ phi\_eq\_psi\_convolution\ (γ\ t\ :\ ℝ)\ (k\ j\ :\ ℕ)\ :}
\leanline{4}{phi\ γ\ (k+j+1)\ t=∫\ s\ in\ 0..t,psi\ γ\ k\ s*psi\ γ\ j\ (t-s)}
\end{leancode}
\leanlabel{How the Lean proof follows the argument}
\texttt{psi\_\allowbreak{}convolution\_\allowbreak{}shift} gives the formal counterpart of this step. The main intermediate facts are \texttt{hd}, \texttt{hc1}, \texttt{hc2}, \texttt{hi} in that order. Named proof applications include \texttt{hi}; the adjacent proof shows their hypotheses and order of use.




\texttt{phi\_\allowbreak{}eq\_\allowbreak{}psi\_\allowbreak{}convolution} proceeds by \texttt{induction j generalizing k with}. The base cases establish the initial claims and the inductive branches use the induction hypotheses. Named proof applications include \texttt{interval\allowbreak{}Integral.\allowbreak{}integral\_\allowbreak{}congr}; the adjacent proof shows their hypotheses and order of use.
\end{leanbox}

\section{\texorpdfstring{The full friction recurrence}{The full friction recurrence}}


Use the just-proved identity with \(j=1\), multiply the integral by \(\gamma\), and substitute \(\gamma\psi_1(t-s)=1-\psi_0(t-s)\). Splitting the scalar integral yields \[\varphi_{k+1}(t)+\gamma\varphi_{k+2}(t)=\psi_{k+1}(t).\] The proof establishes continuity and integrability before splitting; no division by friction is needed. The recurrence therefore includes \(\gamma=0\).



The friction relation is derived by multiplying inside an integral and substituting the division-free balance for the first response. Since no friction denominator is introduced, the identity remains valid at zero friction. This matters for uniform scaling and correction formulas.

\begin{leanbox}{Lean: \texttt{phi\_\allowbreak{}add\_\allowbreak{}friction\_\allowbreak{}phi\_\allowbreak{}succ}}
\leansource{\texttt{phi\_\allowbreak{}add\_\allowbreak{}friction\_\allowbreak{}phi\_\allowbreak{}succ} \quad KineticLangevin/\allowbreak{}ScalarConvolution.lean, lines 40–57}
\begin{leancode}
\leanline{0}{theorem\ phi\_add\_friction\_phi\_succ\ (γ\ t\ :\ ℝ)\ (k\ :\ ℕ)\ :}
\leanline{4}{phi\ γ\ (k+1)\ t+γ*phi\ γ\ (k+2)\ t=psi\ γ\ (k+1)\ t}
\end{leancode}
\leanlabel{How the Lean proof follows the argument}
\texttt{phi\_\allowbreak{}add\_\allowbreak{}friction\_\allowbreak{}phi\_\allowbreak{}succ} gives the formal counterpart of this step. The main intermediate facts are \texttt{hi}, \texttt{hidx}, \texttt{he}, \texttt{hc}, \texttt{h\allowbreak{}I} in that order. The final step is the algebraic normalization or ordered arithmetic displayed in the proof.
\end{leanbox}

\section{\texorpdfstring{Differentiate and integrate the correction coefficients}{Differentiate and integrate the correction coefficients}}


The factored definition \(\varphi_{k+1}(t)=e^{-\gamma t}\int_0^te^{\gamma s}\psi_k(s)\,ds\) has derivative \(\psi_k(t)-\gamma\varphi_{k+1}(t)\). The argument applies the fundamental theorem, product rule, and exact exponential cancellation. At the next index the friction recurrence reduces this derivative to \(\varphi_{k+1}(t)\). Since the next correction vanishes at zero, integrating gives \[\varphi_{k+2}(t)=\int_0^t\varphi_{k+1}(s)\,ds.\] This is the first source identity with its necessary index restriction: the printed assertion at \(k=0\) would incorrectly equate positive \(\varphi_1(t)=te^{-\gamma t}\) with the integral of \(\varphi_0=0\).



The derivative formula has a base-index exception. At positive indices the friction recurrence converts the forcing-minus-damping derivative into the preceding correction response. Checking that exception explicitly prevents an incorrect primitive identity at index zero.

\begin{leanbox}{Lean: \texttt{hasDerivAt\_\allowbreak{}phi\_\allowbreak{}succ}, \texttt{phi\_\allowbreak{}succ\_\allowbreak{}eq\_\allowbreak{}integral\_\allowbreak{}phi}}
\leansource{\texttt{hasDerivAt\_\allowbreak{}phi\_\allowbreak{}succ} \quad KineticLangevin/\allowbreak{}ScalarConvolution.lean, lines 60–74}
\begin{leancode}
\leanline{0}{theorem\ hasDerivAt\_phi\_succ\ (γ\ t\ :\ ℝ)\ (k\ :\ ℕ)\ :}
\leanline{4}{HasDerivAt\ (phi\ γ\ (k+1))\ (psi\ γ\ k\ t-γ*phi\ γ\ (k+1)\ t)\ t}
\end{leancode}
\leansource{\texttt{phi\_\allowbreak{}succ\_\allowbreak{}eq\_\allowbreak{}integral\_\allowbreak{}phi} \quad KineticLangevin/\allowbreak{}ScalarConvolution.lean, lines 77–86}
\begin{leancode}
\leanline{0}{theorem\ phi\_succ\_eq\_integral\_phi\ (γ\ t\ :\ ℝ)\ (k\ :\ ℕ)\ :}
\leanline{4}{phi\ γ\ (k+2)\ t=∫\ s\ in\ 0..t,phi\ γ\ (k+1)\ s}
\end{leancode}
\leanlabel{How the Lean proof follows the argument}
\texttt{has\allowbreak{}Deriv\allowbreak{}At\_\allowbreak{}phi\_\allowbreak{}succ} gives the formal counterpart of this step. The main intermediate facts are \texttt{hc}, \texttt{hd}, \texttt{he} in that order. Its remaining equalities are discharged by the simplifications displayed in the proof.




\texttt{phi\_\allowbreak{}succ\_\allowbreak{}eq\_\allowbreak{}integral\_\allowbreak{}phi} gives the formal counterpart of this step. The main intermediate facts are \texttt{hd}, \texttt{hi} in that order. Named proof applications include \texttt{hi.\allowbreak{}symm}; the adjacent proof shows their hypotheses and order of use.
\end{leanbox}

\end{document}
